% ==================================================================
% id:       lim-001
% topic:    גבולות ורציפות של פונקציות
% subtopic: הגדרת הגבול (ε, δ)
% solution: full
% ==================================================================
\textbf{
בכל אחד מהסעיפים הבאים,
הוכיחו לפי הגדרת הגבול של פונקציה
עם
($\varepsilon$
ו-
$\delta$)
כי:}
\begin{enumerate}\bMs{2}
\item
$\ds{\lim_{x\to 2} \lr{3x^2+5}=17}$
\ifsol
\hint{פרקו את $|f(x)-17|$ למכפלה שאחד מגורמיה הוא $|x-2|$, והגבילו מראש את $\delta$ כדי לחסום את הגורם השני.}
\begin{solbox}
יהי
$\varepsilon>0$. נרצה להראות שקיים $\delta>0$ כך שאם $0<|x-2|<\delta$, אז
\[
.|f(x)-17|=\left|(3x^2+5)-17\right|=3|x-2||x+2|<\varepsilon
\]
נזכור כי
$|x-2|<\delta$
ולכן מאי-שיוויון המשולש נקבל
\[
|x+2|=|x-2+4|\leq |x-2|+4<\delta+4
\]
לכן
$$|f(x)-17|\le 3\cdot \delta\cdot (\delta+4)$$
נדרוש כי
$\delta<1$
ולכן
$3\delta(\delta+4)<15\delta$,
ואז
$$|f(x)-17|<15\delta$$
אז לסיום נבחר
\[
\delta<\min\left\{1,\frac{\varepsilon}{15}\right\}
\]
ואכן, אם $0<|x-2|<\delta$, אז
\[.|f(x)-17|=
\left|(3x^2+5)-17\right|<15\cdot\frac{\varepsilon}{15}=\varepsilon
\]
\end{solbox}
\fi

\item
$\ds{\lim_{x\to-7} |3x+20|=1}$
\ifsol
\hint{השתמשו באי-שוויון המשולש ההפוך $\big||a|-|b|\big|\le|a-b|$.}
\begin{solbox}
יהי $\varepsilon>0$.
לפי אי-שוויון המשולש ההפוך,
שאומר כי
$$||a|-|b||\le |a-b|$$
נקבל כי
\[
.\left||3x+20|-1\right|
=\left||3x+20|-|-1|\right|
\leq |3x+21|=3|x+7|
\]
לכן נבחר
\[
\delta=\frac{\varepsilon}{3}
\]
ואם $0<|x+7|<\delta$, נקבל
\[
.\left||3x+20|-1\right|\leq3|x+7|<\varepsilon
\]
\end{solbox}
\fi

\item
$\ds{\lim_{x\to-1}\frac{3x+5}{2x+4}=1}$
\ifsol
\hint{הביאו את $\frac{3x+5}{2x+4}-1$ למכנה משותף. כדי לחסום את המכנה מלמטה, הגבילו מראש את $\delta$.}
\begin{solbox}
יהי $\varepsilon>0$.
נחפש
$\delta>0$
כך שאם
$|x-(-1)|<\delta$,
אז
מתקיים
\[
.\left|\frac{3x+5}{2x+4}-1\right|
=\left|\frac{x+1}{2x+4}\right|
\]
על המונה יהיה קל לעבוד כי
$|x+1|<\delta$,
והשאלה היא איך נחסום את המכנה.

הרעיון הוא שאם נקח
$x$
מאוד קרוב ל-
$(-1)$
אז
$2x+4$
יהיה מאוד קרוב ל-$2$
וזה מספיק טוב מבחינתנו כי נצליח להקטין את המונה כמה שנרצה (על ידי
$\delta$) אז רק צריך שהמכנה לא ״יעשה בלאגן״. לכן מראש נגביל את
$\delta$
להיות יחסית קטנה ונראה איך זה בפועל חוסם את הערכים שבמכנה:

נדרוש למשל כי
$0<\delta<\frac{1}{2}$
כי אז
$|x+1|<\tfrac{1}{2}$
יגרור
$-\tfrac{3}{2}<x<-\tfrac{1}{2}$,
ולכן
$$.1=2\cdot\left(-\tfrac{3}{2}\right)+4<2x+4<2\cdot\left(-\tfrac{1}{2}\right)+4<3$$

עד כה בחרנו
$\delta<\frac{1}{2}$
ולכן
לכל
$x$
עבורו
$|x+1|<\delta$
מתקיים
$$|f(x)-1|=\left|\frac{x+1}{2x+4}\right|<
\frac{\delta}{|2x+4|}\le \frac{\delta}{1}=\delta$$
כי ראינו שבהכרח
$|2x+4|\ge 1$.
לסיום, נדרוש כי
$\delta<\varepsilon$,
ובסך הכל נדרוש כי
\[
0<\delta<\min\left\{\frac12,\varepsilon\right\}
\]
ואם $0<|x+1|<\delta$, נקבל
\[
|f(x)-1|=\left|\frac{3x+5}{2x+4}-1\right|<\varepsilon
\]
\end{solbox}
\fi
\eMs
\end{enumerate}
